\section*{ЗАКЛЮЧЕНИЕ}
\addcontentsline{toc}{section}{ЗАКЛЮЧЕНИЕ}

В ходе выполнения работы была разработана серверная часть корпоративного мессенджера, обеспечивающая полную независимость от зарубежных поставщиков и санкционных ограничений. Проведённое исследование подтвердило важность автономной backend-системы для российских компаний, гарантирующей стабильную работу коммуникаций и защиту данных на уровне сервера.

Основные результаты работы:

\begin{enumerate}
\item Проведён анализ требований к серверной архитектуре мессенджера, выявлены ключевые риски использования зарубежных решений и определены принципы построения автономной backend-системы.

\item Разработаны концептуальная, логическая и физическая модели данных, оптимизированные для работы серверной части. Спроектирована реляционная схема БД, обеспечивающая эффективное управление пользователями, группами и сообщениями.

\item Реализована серверная логика на Python включающая:

\begin{itemize}
	\item API для взаимодействия с клиентскими приложениями;
	\item модели данных для управления сущностями мессенджера;
	\item механизмы аутентификации и авторизации;
	\item бизнес-логику обработки сообщений;
	\item механизмы отправки/приёма сообщений (личные, групповые, системные);
	\item управление группами пользователей (создание, добавление участников, роли);
	\item поиск по сообщениям и пользователям с фильтрацией;
	\item система аутентификации с валидацией учётных данных.
\end{itemize}

\end{enumerate}

Все требования, объявленные в техническом задании, были полностью реализованы, все задачи, поставленные в начале разработки проекта, были также решены.

Перспективой дальнейшей разработки является расширение функционала мессенджера, добавление поддержки масштабируемых баз данных и внедрение дополнительных механизмов безопасности.